% ============================================================
% Chapter 5: Code Management
% ============================================================
\section{代码管理}

\subsection{如何找出包依赖}

\begin{itemize}[leftmargin=1.5em]
    \item 直接启动试试：\lstinline|application:ensure_all_started(systemd).|
    \item 从 \lstinline|rebar.config| 中查看：\lstinline|{deps, [enough]}.|
    \item 使用命令：\lstinline|rebar3 deps|
    \item 使用命令：\lstinline|rebar3 tree|
    \item 使用命令：\lstinline|make list-deps|
\end{itemize}

\begin{figure}[H]
\centering
\includegraphics[width=0.85\textwidth]{image28}
\caption{rebar3 tree 输出}
\end{figure}

\subsection{如何在 rebar3 工程下设置 application 参数}

使用 \lstinline|config/sys.config| 文件设置 \lstinline|kernel| 等应用的参数。

\subsection{RabbitMQ 工程中修改代码重新编译}

需要删掉整个 \lstinline|ebin| 目录，删掉里面的内容，以及删掉 \lstinline|_build| 都没有用。

\subsection{ez Erlang --- RabbitMQ 的编译方式}

\begin{lstlisting}[style=shellstyle]
make dist DIST_AS_EZS=true
\end{lstlisting}

\subsection{Erlang 批量载入 beam 文件}

\subsection{带有头文件的编译}

参考：\href{https://stackoverflow.com/questions/74700492/erl-make-using-header-file}{Stack Overflow - erl make using header file}

\subsection{Erlang 升级（热代码加载）}

参考：\href{https://medium.com/@kansi/hot-code-loading-with-erlang-and-rebar3-8252af16605b}{Hot Code Loading with Erlang and Rebar3}

\subsection{Loading of Code From Archive Files}

\begin{itemize}[leftmargin=1.5em]
    \item \href{https://www.erlang.org/doc/man/code.html#loading-of-code-from-archive-files}{Erlang -- code}
    \item \href{https://learnyousomeerlang.com/release-is-the-word}{Learn You Some Erlang - Release is the Word}
    \item \href{https://www.erlang.org/doc/man/reltool.html}{Erlang -- reltool}
\end{itemize}
